\section{Reading the formal development}
Chapter III turns three kinds of data into precise finite statements: directed packet operators, coordinate and connection transport, and local obstruction certificates. The preserved finite predecessor is small enough to read as a connected development. Its seven mathematical modules contain the author's constructions and proofs. An aggregate import exposes them together. A separate audit module checks the compiled namespace. The opening chapter describes the finite core and its analytic continuation.

This guide follows the definitions before using their theorem names. A declaration name is a precise destination, but its type determines what has actually been proved. The preserved finite map has 41 rows and 143 distinct targets. Of these, 88 are theorem declarations and 55 are constructions. A construction can be a structure, inductive type, definition, or abbreviation. The count does not mean 143 independent mathematical theorems.

\subsection{Three levels of evidence}
The finite algebraic verification report records a successful aggregate build of 3,076 jobs, the gate that checks parity with the theorem map, and a namespace audit of 468 declarations including 229 entries of kind \code{theoremInfo}. Chapter III also preserves the kernel audit transcript. That transcript contains one axiom report for each of the 143 mapped targets and the final namespace PASS line.

A static reconciliation matches all 41 map rows with their lexical source locations, all 143 explicit \texttt{\#print axioms} directives, and all 143 transcript records. It checked the 16 files in the preserved formal component against the original checksum manifest. The theorem map and transcript hashes also match the values printed in the inherited verification report. These are static checks of the preserved evidence. They do not replace a Lean execution. The opening chapter separately records the compiler and replay runs.

The expanded mathematical text includes additional ordinary mathematical proofs. The finite algebraic formal contract stays fixed. A section of the main text may combine a formally covered finite result with a stronger ordinary theorem. The presence of a related Lean name does not make every statement in that section part of the contract.

\subsection{A route through the source}
\begin{longtable}{>{\raggedright\arraybackslash}p{.29\textwidth}>{\raggedright\arraybackslash}p{.61\textwidth}}
\toprule Module & Reading objective \\ \midrule\endhead
\code{EdgeOperators} & Orientation, broken-word vanishing, and finite walk expansion. \\ \addlinespace
\code{EdgeExtensions} & Endpoint blocks, reachability, and heterogeneous factor types. \\ \addlinespace
\code{Transport} & Exact frames, the minimal connection, and a nontrivial loop. \\ \addlinespace
\code{CoordinateNaturality} & Exact coordinate representation of arbitrary functions. \\ \addlinespace
\code{ReversibleConnections} & Dependent fibers, defect, holonomy, gauges, paths, and blocks. \\ \addlinespace
\code{HolonomyBlocks} & The commutation hypothesis for packetwise gauge covariance. \\ \addlinespace
\code{NonlinearObstructions} & Tree locality and candidate-specific obstruction certificates. \\ \bottomrule
\end{longtable}

The name of a file and the namespace of its declarations need not coincide. For example, the basic packet resolution declarations occur at the end of \code{ReversibleConnections.lean} under \code{Komorebi.PaperIII.HolonomyBlocks}. The complete map gives both the fully qualified name and the file and line. Follow that location when searching a source checkout.

\subsection{Notation that carries mathematical content}
Lean writes a linear map as \code{M →ₗ[R] N} and a linear equivalence as \code{M ≃ₗ[R] N}. The latter includes an inverse and its inverse laws. An arbitrary equivalence is \code{X ≃ Y}. A function \code{X → Y} carries no linear structure. Replacing one of these types with another changes the claim.

Bracketed arguments such as \code{[Semiring R]} are typeclass hypotheses. They supply operations and laws, not merely notation. Braced arguments such as \code{\{R : Type u\}} are implicit parameters inferred from use. \code{Type u} allows the object to live in a universe. Different universe names allow different input, fiber, and coefficient types. \code{autoImplicit false} requires undeclared identifiers to be rejected rather than silently generalized.

\code{Fintype V} supplies finite enumeration. \code{DecidableEq V} supplies a decision procedure for equality. Some definitions use \code{classical} so that a finite sum can be filtered by a proposition without demanding a decision procedure supplied by the user. This logical choice does not turn an abstract graph theorem into an efficient graph algorithm.

Three composition conventions must remain distinct:
\begin{center}
\begin{tabular}{p{.25\textwidth}p{.61\textwidth}}
\toprule Written expression & Execution order \\ \midrule
\code{B * A} & Matrix $A$ acts first, then $B$. \\
\code{g.comp f} & Linear map $f$ acts first, then $g$. \\
\code{e.trans f} & Equivalence $e$ acts first, then $f$. \\ \bottomrule
\end{tabular}
\end{center}
A calculation translated from matrices to \code{LinearEquiv.trans} therefore reverses the written composition pattern. The underlying maps still execute in the same order.

\subsection{What a short proof is doing}
\code{rfl} closes a definitional equality: the two sides reduce to the same expression. \code{ext x} or \code{LinearEquiv.ext} reduces equality of maps to equality at an arbitrary input. \code{simp} uses definitions and registered simplification rules, including inverse laws and zero laws. \code{rw} rewrites by an explicit equality. \code{induction} exposes a finite recursive object, such as a list or tree, and proves the conclusion from the smaller objects.

The printed excerpts are exact source selections, not standalone files. Implicit variables, local instances, namespaces, and earlier lemmas may be outside a selection. Every exact UTF-8 excerpt has a source path, inclusive line range, and hashes in \code{generated/excerpts.json}. Copy the adjacent \code{.lean.txt} file when exact tokens matter. The PDF display uses font fallbacks for otherwise unsupported characters.
